\subsection{Paper Roadmap}

The main goals of this paper are to document, explain, discuss, and evaluate \UT{} and its components.
Some secondary goals of this paper are to be \emph{simple and clear} --- this is why the paper is written in this voice.
When I say ``we'', I am referring to \emph{you and I}.

Our journey is structured thus:
\begin{itemize}
  \item \autoref{sec:proof-of-reflection} constructs the main primitive, \emph{Proof of Reflection}, and covers related topics such as \emph{conversion of work}.
  \item \autoref{sec:constructing-ut} uses \emph{Proof of Reflection} to construct the \emph{Simplex}, \emph{Ultra Terminum} (\UT{}), and its scaling configurations: $\UT{1}$, $\UT{2}$ and $\UT{3}$.
  \item \autoref{sec:practical-considerations} covers the practical considerations for \UT{}: the availability of blocks, state transition, protocol variants, confirmation times, the PoR graph, block-DAGs, DoS attacks, NIPoPoWs, simplex security, and cross-chain transactions.
  \item \autoref{sec:ut-complexity} analyzes the scalability of \UT{}.
  \item \autoref{sec:tiling} constructs a \emph{Tiling of Simplexes} and \UTinf{}, scaling horizontally and unlocking $O(n)$ capacity.
  \item \autoref{sec:attacks} briefly summarizes \UT{}'s resistance to attacks.
  \item \autoref{sec:conclusions} concludes.
\end{itemize}

Additionally, the appendices contain additional protocol analysis and discuss iterations of the simulation (\autoref{sec:por-simulation-experiment}).
